% courses/CDIV2017/Clases/assets/tikzstyles.tex
% ---------------------------------------------------------------------------
% Estilos compartidos para los diagramas del curso Cálculo 1 (CDIV2017).
% Fuente única del "look" visual (colores, grosores, escalas) para que todas
% las figuras (TikZ / pgfplots: conjuntos, rectas numéricas, gráficas de
% funciones) sean consistentes entre clases.
%
% Uso: la clase debe cargar antes `tikz` y `pgfplots` (bloque §6.3 del
%      CLAUDE.md del curso), y luego
%      % courses/CDIV2017/Clases/assets/tikzstyles.tex
% ---------------------------------------------------------------------------
% Estilos compartidos para los diagramas del curso Cálculo 1 (CDIV2017).
% Fuente única del "look" visual (colores, grosores, escalas) para que todas
% las figuras (TikZ / pgfplots: conjuntos, rectas numéricas, gráficas de
% funciones) sean consistentes entre clases.
%
% Uso: la clase debe cargar antes `tikz` y `pgfplots` (bloque §6.3 del
%      CLAUDE.md del curso), y luego
%      % courses/CDIV2017/Clases/assets/tikzstyles.tex
% ---------------------------------------------------------------------------
% Estilos compartidos para los diagramas del curso Cálculo 1 (CDIV2017).
% Fuente única del "look" visual (colores, grosores, escalas) para que todas
% las figuras (TikZ / pgfplots: conjuntos, rectas numéricas, gráficas de
% funciones) sean consistentes entre clases.
%
% Uso: la clase debe cargar antes `tikz` y `pgfplots` (bloque §6.3 del
%      CLAUDE.md del curso), y luego
%      % courses/CDIV2017/Clases/assets/tikzstyles.tex
% ---------------------------------------------------------------------------
% Estilos compartidos para los diagramas del curso Cálculo 1 (CDIV2017).
% Fuente única del "look" visual (colores, grosores, escalas) para que todas
% las figuras (TikZ / pgfplots: conjuntos, rectas numéricas, gráficas de
% funciones) sean consistentes entre clases.
%
% Uso: la clase debe cargar antes `tikz` y `pgfplots` (bloque §6.3 del
%      CLAUDE.md del curso), y luego
%      \input{../assets/tikzstyles.tex}
% (compilar desde el directorio de la clase para que la ruta relativa resuelva).
%
% Paleta propia de este curso (no se reusa la de courses/Fisica3-2015): mismos
% tonos base (azul/rojo/ámbar) para alinear con keybox/notebox/warnbox, pero
% archivo independiente, sin dependencia cruzada entre cursos.
% ---------------------------------------------------------------------------

% ---- Paleta (alineada con las cajas del preámbulo: keybox/notebox/warnbox) ----
\definecolor{figblue}{HTML}{1F4E79}   % azul principal (líneas, keybox, conjuntos)
\definecolor{figred}{HTML}{B03A2E}    % rojo de énfasis (warnbox, región resaltada)
\definecolor{figamber}{HTML}{B8860B}  % ámbar (notebox)
\definecolor{figgray}{HTML}{555555}   % auxiliares, ejes, guías, universo

% ---- TikZ: librerías y estilos reutilizables ----
% Nota: las puntas se declaran como `-{Latex[...]}` y no como `->`. La forma
% con llaves NO contiene un `>` literal, así que es inmune al `>` activo que
% introduce babel español (de todos modos se carga \usetikzlibrary{babel} en
% el bloque §6.3 para cubrir cualquier `->` suelto).
\usetikzlibrary{arrows.meta,calc,angles,quotes,patterns}

\tikzset{
  figline/.style={figblue, line width=0.9pt},
  figaux/.style={figgray, dashed, line width=0.6pt},
  figaxis/.style={figgray, line width=0.7pt, -{Latex[length=2.2mm,width=1.5mm]}},
  % conjuntos (diagramas de Venn / patatoidales)
  figset/.style={figline, fill=figblue!6},
  figsetB/.style={figred, line width=0.9pt, fill=figred!6},
  figsetC/.style={figamber, line width=0.9pt, fill=figamber!10},
  figuniverso/.style={figgray, line width=0.7pt},
  % puntos de la recta / conjuntos numéricos
  figpto/.style={circle, inner sep=0pt, minimum size=5pt, draw=none},
  figptoabierto/.style={circle, inner sep=0pt, minimum size=5pt, draw=figblue, line width=0.9pt, fill=white},
  figptocerrado/.style={circle, inner sep=0pt, minimum size=5pt, draw=figblue, line width=0.9pt, fill=figblue},
  figptoY/.style={figpto, fill=figblue},
  figptoZ/.style={figpto, fill=figred},
  % segmentos gruesos para intervalos sobre la recta
  figintervalo/.style={figblue, line width=2.2pt},
  % vectores/flechas de transformación (traslación, dilatación)
  figvec/.style={figred, line width=1pt, -{Latex[length=2.2mm,width=1.6mm]}},
  % etiquetas
  figlbl/.style={font=\small},
  figlblaux/.style={font=\footnotesize, figgray},
}

% ---- pgfplots: estilo base de las gráficas de funciones ----
\pgfplotsset{
  calcfig/.style={
    width=10cm, height=6cm,
    axis lines=middle,
    xlabel style={font=\small, at={(axis description cs:1,0.5)}, anchor=west},
    ylabel style={font=\small, at={(axis description cs:0.5,1)}, anchor=south},
    tick label style={font=\footnotesize},
    every axis plot/.append style={line width=1.1pt, figblue},
    grid=none,
    legend cell align=left,
    legend style={font=\footnotesize, draw=figgray!40},
    clip=false,
  },
}

% (compilar desde el directorio de la clase para que la ruta relativa resuelva).
%
% Paleta propia de este curso (no se reusa la de courses/Fisica3-2015): mismos
% tonos base (azul/rojo/ámbar) para alinear con keybox/notebox/warnbox, pero
% archivo independiente, sin dependencia cruzada entre cursos.
% ---------------------------------------------------------------------------

% ---- Paleta (alineada con las cajas del preámbulo: keybox/notebox/warnbox) ----
\definecolor{figblue}{HTML}{1F4E79}   % azul principal (líneas, keybox, conjuntos)
\definecolor{figred}{HTML}{B03A2E}    % rojo de énfasis (warnbox, región resaltada)
\definecolor{figamber}{HTML}{B8860B}  % ámbar (notebox)
\definecolor{figgray}{HTML}{555555}   % auxiliares, ejes, guías, universo

% ---- TikZ: librerías y estilos reutilizables ----
% Nota: las puntas se declaran como `-{Latex[...]}` y no como `->`. La forma
% con llaves NO contiene un `>` literal, así que es inmune al `>` activo que
% introduce babel español (de todos modos se carga \usetikzlibrary{babel} en
% el bloque §6.3 para cubrir cualquier `->` suelto).
\usetikzlibrary{arrows.meta,calc,angles,quotes,patterns}

\tikzset{
  figline/.style={figblue, line width=0.9pt},
  figaux/.style={figgray, dashed, line width=0.6pt},
  figaxis/.style={figgray, line width=0.7pt, -{Latex[length=2.2mm,width=1.5mm]}},
  % conjuntos (diagramas de Venn / patatoidales)
  figset/.style={figline, fill=figblue!6},
  figsetB/.style={figred, line width=0.9pt, fill=figred!6},
  figsetC/.style={figamber, line width=0.9pt, fill=figamber!10},
  figuniverso/.style={figgray, line width=0.7pt},
  % puntos de la recta / conjuntos numéricos
  figpto/.style={circle, inner sep=0pt, minimum size=5pt, draw=none},
  figptoabierto/.style={circle, inner sep=0pt, minimum size=5pt, draw=figblue, line width=0.9pt, fill=white},
  figptocerrado/.style={circle, inner sep=0pt, minimum size=5pt, draw=figblue, line width=0.9pt, fill=figblue},
  figptoY/.style={figpto, fill=figblue},
  figptoZ/.style={figpto, fill=figred},
  % segmentos gruesos para intervalos sobre la recta
  figintervalo/.style={figblue, line width=2.2pt},
  % vectores/flechas de transformación (traslación, dilatación)
  figvec/.style={figred, line width=1pt, -{Latex[length=2.2mm,width=1.6mm]}},
  % etiquetas
  figlbl/.style={font=\small},
  figlblaux/.style={font=\footnotesize, figgray},
}

% ---- pgfplots: estilo base de las gráficas de funciones ----
\pgfplotsset{
  calcfig/.style={
    width=10cm, height=6cm,
    axis lines=middle,
    xlabel style={font=\small, at={(axis description cs:1,0.5)}, anchor=west},
    ylabel style={font=\small, at={(axis description cs:0.5,1)}, anchor=south},
    tick label style={font=\footnotesize},
    every axis plot/.append style={line width=1.1pt, figblue},
    grid=none,
    legend cell align=left,
    legend style={font=\footnotesize, draw=figgray!40},
    clip=false,
  },
}

% (compilar desde el directorio de la clase para que la ruta relativa resuelva).
%
% Paleta propia de este curso (no se reusa la de courses/Fisica3-2015): mismos
% tonos base (azul/rojo/ámbar) para alinear con keybox/notebox/warnbox, pero
% archivo independiente, sin dependencia cruzada entre cursos.
% ---------------------------------------------------------------------------

% ---- Paleta (alineada con las cajas del preámbulo: keybox/notebox/warnbox) ----
\definecolor{figblue}{HTML}{1F4E79}   % azul principal (líneas, keybox, conjuntos)
\definecolor{figred}{HTML}{B03A2E}    % rojo de énfasis (warnbox, región resaltada)
\definecolor{figamber}{HTML}{B8860B}  % ámbar (notebox)
\definecolor{figgray}{HTML}{555555}   % auxiliares, ejes, guías, universo

% ---- TikZ: librerías y estilos reutilizables ----
% Nota: las puntas se declaran como `-{Latex[...]}` y no como `->`. La forma
% con llaves NO contiene un `>` literal, así que es inmune al `>` activo que
% introduce babel español (de todos modos se carga \usetikzlibrary{babel} en
% el bloque §6.3 para cubrir cualquier `->` suelto).
\usetikzlibrary{arrows.meta,calc,angles,quotes,patterns}

\tikzset{
  figline/.style={figblue, line width=0.9pt},
  figaux/.style={figgray, dashed, line width=0.6pt},
  figaxis/.style={figgray, line width=0.7pt, -{Latex[length=2.2mm,width=1.5mm]}},
  % conjuntos (diagramas de Venn / patatoidales)
  figset/.style={figline, fill=figblue!6},
  figsetB/.style={figred, line width=0.9pt, fill=figred!6},
  figsetC/.style={figamber, line width=0.9pt, fill=figamber!10},
  figuniverso/.style={figgray, line width=0.7pt},
  % puntos de la recta / conjuntos numéricos
  figpto/.style={circle, inner sep=0pt, minimum size=5pt, draw=none},
  figptoabierto/.style={circle, inner sep=0pt, minimum size=5pt, draw=figblue, line width=0.9pt, fill=white},
  figptocerrado/.style={circle, inner sep=0pt, minimum size=5pt, draw=figblue, line width=0.9pt, fill=figblue},
  figptoY/.style={figpto, fill=figblue},
  figptoZ/.style={figpto, fill=figred},
  % segmentos gruesos para intervalos sobre la recta
  figintervalo/.style={figblue, line width=2.2pt},
  % vectores/flechas de transformación (traslación, dilatación)
  figvec/.style={figred, line width=1pt, -{Latex[length=2.2mm,width=1.6mm]}},
  % etiquetas
  figlbl/.style={font=\small},
  figlblaux/.style={font=\footnotesize, figgray},
}

% ---- pgfplots: estilo base de las gráficas de funciones ----
\pgfplotsset{
  calcfig/.style={
    width=10cm, height=6cm,
    axis lines=middle,
    xlabel style={font=\small, at={(axis description cs:1,0.5)}, anchor=west},
    ylabel style={font=\small, at={(axis description cs:0.5,1)}, anchor=south},
    tick label style={font=\footnotesize},
    every axis plot/.append style={line width=1.1pt, figblue},
    grid=none,
    legend cell align=left,
    legend style={font=\footnotesize, draw=figgray!40},
    clip=false,
  },
}

% (compilar desde el directorio de la clase para que la ruta relativa resuelva).
%
% Paleta propia de este curso (no se reusa la de courses/Fisica3-2015): mismos
% tonos base (azul/rojo/ámbar) para alinear con keybox/notebox/warnbox, pero
% archivo independiente, sin dependencia cruzada entre cursos.
% ---------------------------------------------------------------------------

% ---- Paleta (alineada con las cajas del preámbulo: keybox/notebox/warnbox) ----
\definecolor{figblue}{HTML}{1F4E79}   % azul principal (líneas, keybox, conjuntos)
\definecolor{figred}{HTML}{B03A2E}    % rojo de énfasis (warnbox, región resaltada)
\definecolor{figamber}{HTML}{B8860B}  % ámbar (notebox)
\definecolor{figgray}{HTML}{555555}   % auxiliares, ejes, guías, universo

% ---- TikZ: librerías y estilos reutilizables ----
% Nota: las puntas se declaran como `-{Latex[...]}` y no como `->`. La forma
% con llaves NO contiene un `>` literal, así que es inmune al `>` activo que
% introduce babel español (de todos modos se carga \usetikzlibrary{babel} en
% el bloque §6.3 para cubrir cualquier `->` suelto).
\usetikzlibrary{arrows.meta,calc,angles,quotes,patterns}

\tikzset{
  figline/.style={figblue, line width=0.9pt},
  figaux/.style={figgray, dashed, line width=0.6pt},
  figaxis/.style={figgray, line width=0.7pt, -{Latex[length=2.2mm,width=1.5mm]}},
  % conjuntos (diagramas de Venn / patatoidales)
  figset/.style={figline, fill=figblue!6},
  figsetB/.style={figred, line width=0.9pt, fill=figred!6},
  figsetC/.style={figamber, line width=0.9pt, fill=figamber!10},
  figuniverso/.style={figgray, line width=0.7pt},
  % puntos de la recta / conjuntos numéricos
  figpto/.style={circle, inner sep=0pt, minimum size=5pt, draw=none},
  figptoabierto/.style={circle, inner sep=0pt, minimum size=5pt, draw=figblue, line width=0.9pt, fill=white},
  figptocerrado/.style={circle, inner sep=0pt, minimum size=5pt, draw=figblue, line width=0.9pt, fill=figblue},
  figptoY/.style={figpto, fill=figblue},
  figptoZ/.style={figpto, fill=figred},
  % segmentos gruesos para intervalos sobre la recta
  figintervalo/.style={figblue, line width=2.2pt},
  % vectores/flechas de transformación (traslación, dilatación)
  figvec/.style={figred, line width=1pt, -{Latex[length=2.2mm,width=1.6mm]}},
  % etiquetas
  figlbl/.style={font=\small},
  figlblaux/.style={font=\footnotesize, figgray},
}

% ---- pgfplots: estilo base de las gráficas de funciones ----
\pgfplotsset{
  calcfig/.style={
    width=10cm, height=6cm,
    axis lines=middle,
    xlabel style={font=\small, at={(axis description cs:1,0.5)}, anchor=west},
    ylabel style={font=\small, at={(axis description cs:0.5,1)}, anchor=south},
    tick label style={font=\footnotesize},
    every axis plot/.append style={line width=1.1pt, figblue},
    grid=none,
    legend cell align=left,
    legend style={font=\footnotesize, draw=figgray!40},
    clip=false,
  },
}
